\section{Interfacial Spin Configurations and the Microscopic Origin of Exchange Bias at the IrMn/CoFeB Interface \href{https://doi.org/10.1016/j.fmre.2026.01.003}{[\textit{Fundamental Research, 8 Jan (2026)}]}}

\section*{Main points from the Paper}
\subsection*{\centering\underline{Abstract}}
\hspace{3cm} Exchange bias fields at antiferromagnet/ferromagnet (AFM/FM) interfaces' physical origin remains incompletely understood. While the most of the spins are strongly coupled to the adjacent CoFeB layer, a small fraction remains pinned to the underlying IrMn underlayer. Micromagnetic simulations indicates that an imbalance in the quantity of antiparallel pinned spins contributes to the observed variation in exchange bias.

\subsection*{\centering\underline{I. Introduction}}
\begin{enumerate}
    \item Study of AFM/FM interface has greater impact in the study of spintronics. Exchange bias, being a notably in the AFM/FM structures have a significant in the performance of spintronic devices.
    \item \hl{Exchange bias (EB)} refers to the shift of the magnetization hysteresis loop caused by interfacial coupling between the AFM and FM layers.
    \item EB was generally used to stabilize the magnetization of the pinnned or reference layer in spin-valve structures and MTJs. Moreover, spin-orbit torque (SOT) switching of perpendicular magnetization without an external magnetic field can be realized by introducing an in-plane EB field.
    \item In this work CoFeB/MgO structure provides strong perpendicular magnetic anisotropy and large tunneling magnetoresistance.
    \item Due to the recent development in the electrical control of EB via SOT has been demonstrated, enabling the development of a novel storage prototype known as AFM random-access memory.
    \item Earlier studies  have explained that EB originates from fully uncompensated AFM spins [\textcolor{red}{ref. 39}], planar domain-wall model to reconcile the quantitative discrepancies between theoretical predictions and experimentally measured the EB fields [\textcolor{red}{ref. 40}]. Subsequently, Malozemoff developed the random-field model to account for the influence of interface roughness and defects [\textcolor{red}{ref. 42}].
    \item These established that the coupling between FM spins and pinned AFM  spins is a key microscopic mechanism responsible for the EB field.
    \item Recently, the electrical control of EB via SOT have shown that spin current can partially switch the EB field, this is particularly attributed to the formation of AFM multidomains during manipulation process, suggesting that the orientations of interfacial AFM spins, and their contributions to the exchange bias field, are more complex than previously studied.
    \item X-ray magnetic circular dichroism (XMCD)  is employed as a high-precession technique for probing and analyzing AFM/FM interfaces.
    \item Measurements are taken as;
    \begin{itemize}
        \item XMCD measurements at different temperature is done.
        \item Element-specific XMCD hysteresis loops are measured to determine the orientation of AFM pinned spins in samples annealed under different magnetic fields.
    \end{itemize}
    \item Micromagnetic simulations show that variations in the exchange bias field are primarily governed by the quantitative imbalance between these pinned-spin populations.
\end{enumerate}

\ding{107} \hl{Materials and Methods section is not included in these notes}.
\subsection*{\centering\underline{II. Results and Discussion}}
\begin{enumerate}
    \item The multilayer exhibits perpendicular magnetic anisotropy along with an out-of-plane exchange bias field. They have observed the monotonically increment in the magnitude of the EB with the annealing field and saturates at $\pm1$ T.
    \item 
\end{enumerate}